\appendix

%

\section{Birg{\'e}'s algorithm as a semi-agnostic learner}
\label{sec:birge-agnostic}

In this section we briefly explain why Birg{\'e}'s algorithm~\cite{Birge:87b}
also works in the semi-agnostic setting,
{thus justifying the claims about its performance made in the
statement of \expref{Theorem}{thm:birge-monotone}.}
To do this, we need to explain
his approach. For this, we will need the following {fact (which follows as a special case of the VC inequality, {\expref{Theorem}{thm:vc-inequality}}}), 
which gives
a tight bound on the number of samples required to
learn an arbitrary distribution with respect to \emph{total variation distance}.

\begin{fact} \label{thm:folklore}
Let $p$ be any distribution over $[n]$. We have: $\E [ \dtv ( p, \wh{p}_m) ]  =  O( \sqrt{n/m})$.
\end{fact}


Let $p$ be a non-increasing distribution over $[n]$.
(The analysis for the non-decreasing case is identical.)
Conceptually, we view algorithm $\ANI$ as working in three steps:

\begin{itemize}
\item  In the first step, it partitions the set $[n]$ into a carefully chosen
set $I_1, \ldots, I_{\ell}$ of consecutive intervals, with $\ell = O(m^{1/3} \cdot (\log n)^{2/3})$.
Consider the \emph{flattened} distribution $p_f$ over $[n]$ obtained from $p$ by averaging the weight
that $p$ assigns to each interval over the entire interval. That is, for $j \in [\ell]$ and
$i \in I_j$, $p_f(i) = \sum_{t \in I_j} p(t) / |I_j|$.
Then a simple argument given in~\cite{Birge:87b} gives
that 
\[
\dtv(p_f, p) = O\left( ({\log n / (m+1)})^{1/3} \right)\,.
\]


\item Let $p_r$ be the \emph{reduced} distribution corresponding to $p$ and the partition $I_1, \ldots, I_{\ell}$.
That is, $p_r$ is a distribution over $[\ell]$ with $p_r(i) = p(I_i)$ for $i \in [\ell]$. In the second step, the algorithm uses the $m$ samples to learn $p_r$. (Note that $p_r$ is not necessarily monotone.)
After $m$ samples, one obtains a hypothesis $\widehat{p_r}$ such that
\[
\E [\dtv(p_r, \widehat{p_r})] =  O\left( \sqrt{\ell / m} \right) =
O\left( ({\log n / (m+1)})^{1/3} \right)\,.
\]
The first equality follows from \expref{Fact}{thm:folklore} (since $p_r$ is distribution over $\ell$ elements) and
the second inequality follows from the choice of $\ell$.

\item Finally, the algorithm outputs the flattened hypothesis $(\widehat{p_r})_f$ over $[n]$
corresponding to  $\widehat{p_r}$, \ie, obtained by $\widehat{p_r}$ by subdividing
the mass of each interval uniformly within the interval. It follows from the above two steps that
\[
\E[\dtv ( (\widehat{p_r})_f, p_f)] =O\left( ({\log n / (m+1)})^{1/3}
\right)\,.
\]

\item The combination of the first and third steps yields that
\[
\E[\dtv ( (\widehat{p_r})_f, p)] =O\left( ({\log n / (m+1)})^{1/3}
\right)\,.
\]


\end{itemize}

\medskip

The above arguments are entirely due to Birg{\'e}~\cite{Birge:87b}.
We now explain how his analysis can be extended to show that his algorithm is
in fact a semi-agnostic learner as claimed in \expref{Theorem}{thm:birge-monotone}.
To avoid clutter in the expressions below let us fix
\[
\delta: =O\left( ({\log n / (m+1)})^{1/3} \right)\,.
\]

The second and third steps in the algorithm description
above are used to learn the distribution $p_f$ to variation distance $\delta$.
Note that these steps do not use the assumption that $p$ is non-increasing.
The following claim, which generalizes Step 1 above, says that if $p$ is $\tau$-close to non-increasing,
the flattened distribution $p_f$ (defined as above) is $(2\tau+\delta)$-close to $p$.
Therefore, it follows that, for such a distribution $p$,
algorithm $\ANI$ succeeds with expected (total variation distance) error
$(2\tau+\delta)+\delta$.

\noindent We have:

\begin{claim}
Let $p$ be a distribution over $[n]$ that is $\tau$-close to non-increasing.
Then, the flattened distribution $p_f$ (obtained from $p$ by averaging its
weight on every interval $I_j$) satisfies
\[
\dtv(p_f,p) \leq (2\tau+\delta)\,.
\]
\end{claim}

\begin{proof}
Let $p^{\downarrow}$ be the non-increasing distribution that is $\tau$-close to $p$.
Let $\tau_j$ denote the $L_1$-distance between $p$ and $p^{\downarrow}$ in the interval $I_j$.
Then, we have that
\begin{equation} \label{eqn:one}
\sum_{j=1}^{\ell} \tau_j \leq \tau\,.
\end{equation}

By Birg{\'e}'s arguments, it follows that the flattened distribution $(p^{\downarrow})_f$
corresponding to $p^{\downarrow}$ is $\delta$-close to $p^{\downarrow}$,
hence $(\tau+\delta)$-close to $p$. That is,
\begin{equation} \label{eqn:two}
\dtv \left( (p^{\downarrow})_f, p \right)  \leq \tau+\delta\,.
\end{equation}
We want to show that
\begin{equation} \label{eqn:three}
\dtv \left( (p^{\downarrow})_f, p_f \right) \leq \tau\,.
\end{equation}
Assuming (\ref{eqn:three}) holds, we can conclude by the triangle inequality that
$$\dtv \left(p, p_f \right) \leq 2\tau+\delta$$
as desired.

Observe that, by assumption, $p$ and $p^{\downarrow}$ have
$L_1$-distance at most $\tau_j$ in each $I_j$ interval.
In particular, this implies that, for all $j \in [\ell]$, it holds
$$ \left| p(I_j) - p^{\downarrow}(I_j) \right| \leq \tau_j\,.$$
Now note that, within each interval $I_j$, $p_f$ and $(p^{\downarrow})_f$ are both uniform.
Hence, the contribution of $I_j$ to the variation distance between  $p_f$ and $(p^{\downarrow})_f$
is at most $|p(I_j) - p^{\downarrow}(I_j)|$.

Therefore, by (\ref{eqn:one}) we deduce
$$\dtv (p_f , (p^{\downarrow})_f ) \leq \tau$$
which completes the proof of the claim.
\end{proof}

\section{Hypothesis testing} \label{sec:choosehypothesis}

Our hypothesis testing routine \algorithml{Choose-Hypothesis}$^p$ runs a simple ``competition'' to choose a winner between two candidate hypothesis distributions $h_1$ and $h_2$  over $[n]$ that it is given in the input either explicitly, or in some succinct way.  We show that if at least one of the two candidate hypotheses is close to the target distribution $p$, then with high probability over the samples drawn from $p$ the routine selects as winner a candidate that is close to $p$.
  This basic approach of running a competition between candidate hypotheses is quite similar to the ``Scheff\'e estimate'' proposed by Devroye and Lugosi (see~\cite{DL96,DL97}
and Chapter~6 of~\cite{DL:01}), which in turn built closely on the work of~\cite{Yatracos85}, but
there are some small differences between our approach and theirs; the~\cite{DL:01} approach uses a notion of the ``competition'' between two hypotheses which is not symmetric under swapping the two competing hypotheses, whereas our competition is {symmetric}.
%

We now prove \expref{Theorem}{thm:choosehypothesis}.

\begin{proof}[Proof of \expref{Theorem}{thm:choosehypothesis}]
Let ${\cal W}$ be the support of $p$. To set up the competition between $h_1$ and $h_2$, we define the following subset of ${\cal W}$:
\begin{align}
&{\cal W}_1={\cal W}_1(h_1,h_2) := \left\{w \in \mathcal{W}~\vline~h_1(w) > h_2(w) \right\}\,. \label{eq:W1}
%
%
\end{align}

\noindent Let then $p_1 = h_1({\cal W}_1)$ and $q_1 = h_2({\cal W}_1)$. Clearly, $p_1 > q_1$ and
 $\dtv(h_1, h_2) = p_1-q_1$.

%

The competition between $h_1$ and $h_2$ is carried out as follows:

\begin{enumerate}
\item  If $p_1-q_1\leq 5 \eps'$, declare a draw and return either $h_i$. Otherwise:

\item  Draw 
  \[
  m=O\left(\frac{\log(1/\delta')}{\eps'^2}\right)
  \]
  samples $s_1,\ldots,s_m$ from $p$, and let 
  \[
  \tau = \frac{1}{m}\, \bigl| \{i~|~s_i \in {\cal W}_1 \}\bigr|
  \]
  be the fraction of  samples that fall inside ${\cal W}_1$.


\item  If $\tau > p_1- \frac{3}{2} \eps'$, declare $h_1$ as winner and return $h_1$; otherwise,

\item  if $\tau < q_1+ \frac{3}{2} \eps'$, declare $h_2$ as winner and return $h_2$; otherwise,

\item  declare a draw and return either $h_i$.

\end{enumerate}

It is not hard to check that the outcome of the competition does not depend on the ordering of the pair of distributions provided in the input; that is, on inputs $(h_1,h_2)$ and $(h_2,h_1)$ the competition outputs the same result for a fixed sequence of samples $s_1,\ldots,s_m$ drawn from $p$.

The correctness of \algorithml{Choose-Hypothesis} is an immediate consequence of the following lemma.

%

\begin{lemma} \label{lem:kostas3}
Suppose that $\dtv(p,h_1) \leq \eps'$. Then:
\begin{itemize}

\item[(i)] If $\dtv(p, h_2)>6 \eps'$, then the probability that the competition between
$h_1$ and $h_2$ does not declare $h_1$ as the winner is at most $e^{- {m \eps'^2 /2 } }$.  (Intuitively,
if $h_2$ is very bad then it is very likely that $h_1$ will be declared winner.)

\item[(ii)]  If $\dtv(p, h_2)>4 \eps'$, the probability that the competition between $h_1$
and $h_2$ declares $h_2$ as the winner is at most $e^{- {m \eps'^2/ 2 } }$.  (Intuitively, if $h_2$
is only moderately bad then a draw is possible but it is very unlikely that $h_2$ will be declared winner.)
\end{itemize}
\end{lemma}


%
\begin{proof}
Let $r=p({\cal W}_1)$. The definition of the total variation distance
implies that $|r-p_1| \le \eps'$. Let us define the $0/1$ (indicator)
random variables $\{Z_j\}_{j=1}^m$ as $Z_j=1$ iff $s_j \in {\cal
  W}_1$. Clearly, \[
\tau=\frac{1}{m} \sum_{j=1}^m Z_j\qquad\text{and}\qquad
\mathbb{E}[\tau]=\mathbb{E}[Z_j]=r\,.
\]
Since the $Z_j$'s are mutually independent, it follows from the
Chernoff bound that 
\[
\Pr[\tau \le r-{\eps'/2}] \le e^{- {m \eps'^2/2 } }\,.
\]
Using $|r-p_1| \le \eps'$ we get that 
\[
\Pr[\tau\le p_1-{3\eps'/2}] \le e^{- {m \eps'^2/2 } }\,.
\]
\begin{itemize}
\item For part (i): If $\dtv(p, h_2) > 6 \eps'$, from the triangle inequality we get that $p_1-q_1=\dtv(h_1, h_2) > 5 \eps' $.  Hence, the algorithm will go beyond Step 1, and with probability at least $1-e^{- {m \eps'^2/2 } }$, it will stop at Step 3, declaring $h_1$ as the winner of the competition between $h_1$ and $h_2$.

\item For part (ii):  If $p_1-q_1 \leq 5 \eps'$ then the competition declares a draw, hence $h_2$ is not the winner. Otherwise we have $p_1 - q_1 > 5\eps'$ and the above arguments imply that the competition between $h_1$ and $h_2$ will declare $h_2$ as the winner with probability at most $e^{- {m \eps'^2/2 } }$.
\end{itemize}
This concludes the proof of \expref{Lemma}{lem:kostas3}.
\end{proof}

\noindent
The proof of the theorem is now complete.
\end{proof}

\section{Using the hypothesis tester} \label{sec:ht}

In this section, we explain in detail how we use
the hypothesis testing algorithm \algorithml{Choose-Hypothesis}
throughout this paper. In particular, the algorithm
\algorithml{Choose-Hypothesis} is used in the following places:
\begin{itemize}

\item In Step 4 of algorithm \hyperref[alg:learn-kmodal-simple]{\algorithml{Learn-kmodal-simple}} we need
an algorithm $\ANI_{\delta'}$ (resp. $\AND_{\delta'}$) that learns
a non-increasing (resp.\ non-increasing) distribution within total variation distance
$\eps$ and confidence $\delta'$. Note that the corresponding algorithms
$\ANI$ and $\AND$ provided by \expref{Theorem}{thm:birge-monotone} have confidence $9/10$.
To boost the confidence of $\ANI$ (resp. $\AND$)
we run the algorithm $O(\log(1/\delta'))$ times and use \algorithml{Choose-{\allowbreak}Hypothesis}
in an appropriate tournament procedure to select among the candidate hypothesis distributions.

\item In Step 5 of algorithm \hyperref[alg:learn-kmodal-simple]{\algorithml{Learn-kmodal-simple}} we need to
select among two candidate hypothesis distributions (with the
promise that at least one of them is close to the true
conditional distribution). In this case, we run \algorithml{Choose-Hypothesis} once to select between
the two candidates.

\item Also note that both algorithms \hyperref[alg:learn-kmodal-simple]{\algorithml{Learn-kmodal-simple}} and \hyperref[alg:learn-kmodal]{\algorithml{Learn-kmodal}} generate
an $\eps$-accurate hypothesis with probability $9/10$. We would like to boost the probability
of success to $1-\delta$. To achieve this we again run the corresponding algorithm $O(\log(1/\delta))$
times and use \algorithml{Choose-{\allowbreak}Hypothesis} in an appropriate tournament
to select among the candidate hypothesis distributions.
\end{itemize}

We now formally describe the ``tournament'' algorithm to boost the confidence to $1-\delta$.

\begin{lemma} \label{lem:log-cover-size}
Let $p$ be any distribution over a finite set $\mathcal{W}$.
Suppose that ${\cal D}_\eps$ is a collection of $N$ distributions over $\mathcal{W}$ such that
there exists $q \in {\cal D}_\eps$ with $\dtv(p,q) \leq \eps$.
Then there is an algorithm that uses $O(\eps^{-2}\log N \log(1/\delta))$ samples from
$p$ and with probability $1-\delta$ outputs a distribution $p' \in {\cal D}_{\eps}$ that satisfies
$\dtv(p,p') \leq 6\eps$.
\end{lemma}


Devroye and Lugosi (Chapter 7 of~\cite{DL:01}) prove a similar result by having all pairs
of distributions in the cover compete against each other using their notion of a competition, but again there
are some small differences:  their approach chooses a distribution in the cover which wins the
maximum number of competitions, whereas our algorithm chooses
a distribution that is never defeated (\ie, won or achieved a draw against all other distributions
in the cover).  Instead we follow the approach from~\cite{DDS12stoclong}.

\medskip

\begin{proof}
The algorithm %
performs a tournament by running the competition
\algorithml{Choose-Hypothesis}$^p(h_i,h_j,\eps,$ $\delta/(2N))$ for every pair of distinct distributions
$h_i,h_j$ in the collection ${\cal D}_\eps$.  It outputs a distribution $q^{\star} \in {\cal D}_\eps$
that was never a loser (\ie, won or achieved a draw in all its competitions).
If no such distribution exists in ${\cal D}_\eps$ then the algorithm outputs
``failure.''

By definition, there exists some $q \in {\cal D}_\eps$ such that $\dtv(p,q) \leq \eps$.
We first argue that with high probability this distribution $q$ never
loses a competition against any other $q' \in {\cal D}_\eps$
(so the algorithm does not output ``failure'').  Consider any $q'
\in {\cal D}_\eps$. If $\dtv(p, q') > 4 \eps$, by \expref{Lemma}{lem:kostas3}(ii) the probability that $q$ loses to
$q'$ is at most $2e^{-m \eps^2 /2} = O(1/N)$. On the other hand, if $\dtv(p, q') \leq 4 \delta$, the triangle
inequality gives that $\dtv(q,q') \leq 5 \eps$ and thus $q$ draws against $q'$.  A union
bound over all $N$ distributions in ${\cal D}_{\eps}$ shows that with probability $1-\delta/2$, the distribution $q$ never loses a competition.

We next argue that with probability at least $1-\delta/2$, every distribution $q' \in {\cal D}_\eps$ that never loses has small variation distance from $p$. Fix a distribution $q'$ such that $\dtv(q',p) > 6 \eps$; \expref{Lemma}{lem:kostas3}(i) implies that $q'$ loses to $q$ with probability $1 - 2e^{-m \eps^2 /2} \geq 1 - \delta/(2N)$.  A union bound gives that with probability $1-\delta/2$, every distribution $q'$ that has $\dtv(q', p) > 6 \eps$ loses some competition.

Thus, with overall probability at least $1-\delta$, the tournament does not output ``failure'' and outputs some distribution $q^{\star}$ such that $\dtv(p, q^{\star})$ is at most $6 \eps$.
This proves the lemma.
\end{proof}


We now explain how the above lemma is used in our context: Suppose we perform $O(\log(1/\delta))$ runs of a learning algorithm that constructs an $\eps$-accurate hypothesis with probability at least $9/10$.  Then, with failure probability at most $\delta/2$, at least one of the hypotheses generated is $\eps$-close to the true distribution in variation distance. Conditioning on this good event, we have a collection of distributions with cardinality $O(\log(1/\delta))$ that satisfies the assumption of the lemma. Hence, using
$O\left((1/\eps^2) \cdot \log\log(1/\delta) \cdot \log(1/\delta)\right) $ samples we can learn to accuracy
$6\eps$ and confidence $1-\delta/2$. The overall sample complexity is $O(\log(1/\delta))$ times the sample complexity of the 
{learning algorithm run with confidence $9/10$,} 
plus this additional $O\left((1/\eps^2) \cdot \log\log(1/\delta) \cdot \log(1/\delta)\right)$ term.

In terms of running time,we make the following easily verifiable remarks: When the hypothesis testing algorithm \algorithml{Choose-Hypothesis}  is run on a pair of distributions that are produced by Birg{\'e}'s algorithm, its running time is polynomial in the succinct description of these distributions, \ie, in $\log^2(n) / \eps$. Similarly, when \algorithml{Choose-Hypothesis} is run on a pair of outputs of \hyperref[alg:learn-kmodal-simple]{\algorithml{Learn-kmodal-simple}} or \hyperref[alg:learn-kmodal]{\algorithml{Learn-kmodal}}, its running time is polynomial in the succinct description of these distributions. More specifically, in the former case, the succinct description has bit complexity $O \left( k \cdot \log^2 (n) /\eps^2 \right)$ (since the output consists of $O(k/\eps)$ monotone intervals, and the conditional distribution on each interval is the output of Birg{\'e}'s algorithm for that interval). In the latter case, the succinct description has bit complexity
$O\left( k \cdot \log^2 (n) /\eps \right)$, since the algorithm \hyperref[alg:learn-kmodal]{\algorithml{Learn-kmodal}} constructs  only $k$ monotone intervals. Hence, in both cases, each execution of the testing algorithm performs  $\poly(k, \log n , 1/\eps)$ bit operations. Since the tournament invokes the algorithm \algorithml{Choose-Hypothesis} $O(\log^2(1/\delta))$ times (for every pair of distributions in our pool of $O(\log(1/\delta))$ candidates) the upper bound on the running time follows.
